\section{Discussion: Evidence Strength and Open Questions}

The strongest direct evidence concerns the phase transition of BaZn$_2$Si$_2$O$_7$, the Ba\slash Sr solid solution, and Zn-site substitution, where HT\nobreakdash-XRD, single-crystal XRD, and dilatometry corroborate one another \cite{lin1999phase,thieme2015ba1,kerstan2012thermal,thieme2016new}. On the glass-crystallization side, nucleating-agent and microstructural effects have been studied systematically, but glass composition, heating rate, and grain size vary from study to study, so ``phase stability'' and ``macroscopic thermal expansion of the sample'' cannot be treated as equivalent \cite{thieme2016thermal,vladislavova2018crystallisation}.

The first major gap is the ternary BaO--ZnO--SiO$_2$ phase diagram and its thermodynamic functions. For now, phase fields can only be inferred by piecing together binary phase equilibria, laboratory XRD, and thermal analysis, and theory indicates that vibrational free energy and temperature effects must be included \cite{erlebach2021phase}. The second gap is direct structural evidence for Ba$_5$Y$_{12}$Zn[O(SiO$_4$)]$_8$. That Ba, Y, Zn, and silicate groups coexist in one formula does not establish an isostructural relationship; verifiable diffraction data and coordination refinement are needed as well.

We recommend a safe, falsifiable experimental path: prepare composition-gradient samples on a small scale in a closed high-temperature furnace with vetted crucible materials; collect room- and high-temperature powder XRD, TG\nobreakdash-DSC, Raman spectra, and microscale compositional maps in parallel; and reserve synchrotron or single-crystal structure determination for samples that show a single diffraction phase. This path follows the complementary use of HT\nobreakdash-XRD and dilatometry in the BaZn$_2$Si$_2$O$_7$ studies \cite{kerstan2012thermal,thieme2016very}; unverified Y-rich formulations are not presented as ready-to-run protocols.

Structural descriptors can help narrow the candidate space. Negative-thermal-expansion studies show that tetrahedral rotation, chain flexibility, and coordination-polyhedron voids serve as useful screening features; descriptors, however, only generate candidates and cannot replace phase-equilibrium experiments \cite{kireeva2022pauling,erlebach2021phase}.
